\documentclass[acmsmall,screen,review,anonymous]{acmart}
% Review manuscript based on the ACM acmsmall conference sample.
\acmJournal{PACMSE}
\acmConference[ISSTA '27]{ACM SIGSOFT International Symposium on Software Testing and Analysis}{September 7--10, 2027}{Singapore}
\setcopyright{none}
\acmDOI{}
\acmYear{2027}
\copyrightyear{2027}
\acmISBN{}
\title{What Does a Passing Regression Test Tell Us About a Solidity Repair?}
\author{Anonymous Author(s)}
\renewcommand{\shortauthors}{Anonymous Author(s)}
\begin{document}

\begin{abstract}
A maintainer who accepts a patch because it passes the regression test in an audit report still has to decide whether other callers can rely on the repaired contract. In an archived Solidity repair study, 18 of 96 Angstrom patches that passed the supplied example failed another frozen security check. In a separate Ekubo study, follow-up checks found the reported defect in 17 of 59 patches that had passed both the original checks and the ordinary-use suite. Neither study varied which example of the same defect the agent received under random assignment. We describe a benchmark that separates three questions: whether a patch corrects the reported behaviour while preserving specified ordinary use, how often the same attempts also satisfy a distinct withheld behaviour on a qualified subset, and whether choosing which of two examples to show changes what the agent repairs. Among 160 external findings reviewed for scope, 38 cleared documentary review, 37 built on both historical revisions, and 18 have evidence of affected ordinary use on both revisions offline. Five cases have evidence pages ready for human admission. The planned repair experiment awaits a signed roster. This first version reports the development studies and case-construction results; the comparative repair outcomes remain open.
\end{abstract}

\ccsdesc[500]{Software and its engineering~Software testing and debugging}
\ccsdesc[300]{Security and privacy~Software security engineering}
\keywords{automated program repair, Solidity, regression testing, repair evaluation}
\maketitle

\section{Introduction}
A maintainer receives a patch for a price-history error and reruns the test from the repair ticket. The test passes. Another caller can read the same pool's history through a different contract interface. Should the maintainer accept the patch for that caller too? This was a concrete distinction in an earlier study of Ekubo. Some patches returned the expected history through the library interface while the direct contract interface still returned an incorrect answer. One passing test settled the question it asked. The maintainer's decision depended on a second question.

Solidity makes the distinction observable because different callers can depend on the same persistent contract state. A change that prevents a reported failure may also interrupt a legitimate operation. Conversely, a patch can keep ordinary workflows running while leaving the audited defect visible under another use. EVMPatch made both aspects measurable by checking the effect of Ethereum patches against harmful and legitimate historical transactions~\cite{rodler2021evmpatch}. For a repair agent that reads a regression example in the ticket, we want to distinguish learning what the contract should do from fitting what that example happens to exercise.

Development records give reasons to expect both incomplete repairs and ceilings. On one finding, a second check rejected patches that had passed the supplied example. On another, new checks rejected patches that had already passed a battery of checks and the ordinary-use suite. A different configuration on the latter finding passed even the expanded battery across all its counted sessions. These observations establish that the choice of measurement matters on these cases. They leave open whether changing the example given to a repair agent would change the patch.

The missing comparison is specific. For the same audited defect, keep the ticket, code, working environment and grader fixed. Vary only which of two legitimate examples a maintainer includes. Under either assignment, ask separately whether each behaviour now works. A finding can contribute to a benchmark of reported repairs even when its public history supports only one observable obligation. Counting a second copy of that same obligation as evidence of transfer would make the apparent breadth of a repair depend on how many equivalent tests a curator wrote.

We therefore ask three questions with explicit denominators. \textbf{RQ1 (reported behaviour).} On admitted findings, how often does a repair configuration deliver a \emph{working repair of the reported behaviour}? The submitted patch builds, preserves the specified ordinary use, and corrects the defect shown in example A for a reason supported by the recorded assessment. \textbf{RQ2 (another use).} Among all assigned A attempts on findings with a qualified hidden behaviour H, what fraction satisfies both A and H while preserving ordinary use, and how much lower is that fraction than the primary rate on those same attempts? \textbf{RQ3 (example choice).} Among findings with two legitimate alternatives A and B and a common withheld H, does randomly showing A or B preferentially improve the behaviour illustrated? RQ1 and RQ2 use standard reports showing A. RQ3 adds the alternative assignment.

This manuscript presents two contributions in hand and a prospective experiment:
\begin{itemize}
\item A comparison of the evidence behind three archived development results, including their different denominators and when the graders were constructed.
\item A case-admission and measurement protocol, with an account of which public findings have reached each qualification step and which evidence is still missing.
\item A registered primary outcome and a design for comparing repair configurations and the identity of the supplied example. The performance estimates are due after admission, assignment and grading.
\end{itemize}

\section{What earlier repair studies measure}
\paragraph{Changing the tests available during repair.}
Smith et al.\ evaluated patch overfitting using held-out test suites~\cite{smith2015overfitting}. Motwani et al.\ came closer to our intervention: within defects, they sampled failing tests while holding passing tests fixed and examined patch quality as the number of failing tests changed~\cite{motwani2022quality}. Modern agents also use tests as feedback. Ahmed et al. studied test overfitting on SWE-bench, including access to hidden developer tests~\cite{ahmed2026overfitting}. TDFlow contrasts developer-supplied tests with tests produced by the system~\cite{han2026tdflow}. These studies motivate controlling the amount and origin of test information. We fix the number of reported examples at one and estimate two directional contrasts within each defect. This lets an intrinsically easier behaviour succeed often under both assignments without being mistaken for an advantage from showing it.

\paragraph{Measuring a security repair.}
Vul4J links real Java vulnerabilities to historical revisions and executable evidence~\cite{bui2022vul4j}. SEC-bench evaluates end-to-end agents on real security tasks~\cite{lee2025secbench}. EVMbench already evaluates smart-contract patching using existing project tests and hidden checks~\cite{wang2026evmbench}. The PatchBench preprint examines how security and semantic checks change patch assessments in C and C++~\cite{shen2026patchbench}. Independent testing, Solidity patching, and failure under stronger grading thus have direct precedents. The proposed contribution concerns \emph{which reported behaviour to illustrate} and the conditions under which a second observation can support an inference about another use. We limit this distinction to the studies compared here.

\section{Evidence from two development findings}
\label{sec:development}
The archived studies provide examples of what a grader can miss and when extra grading stops distinguishing patches. Their unit of observation is a repair attempt on a development finding. Table~\ref{tab:development} gives the eligibility condition for each denominator. The denominator changes across rows because the underlying studies asked different questions.

\begin{table}[t]
\caption{Development observations. Each denominator is conditional on the entry rule in the middle column. Each row concerns one finding.}
\label{tab:development}
\small
\begin{tabular}{@{}p{.18\linewidth}p{.38\linewidth}p{.34\linewidth}@{}}
\toprule
Finding and setting & Which attempts enter & What the additional evidence showed \\
\midrule
Angstrom, original study & 96 patches passing the supplied example among 108 sessions shown one & 18 of 96 failed at least one other frozen security check. \\
Ekubo, GPT-5.4-mini & 59 of 144 submissions passed the original battery and ordinary-use suite & Later checks found the audited defect in 17 of 59. \\
Ekubo, Claude Code with Opus & 48 counted sessions across four conditions & All 48 passed the original, ordinary-use and later checks mechanically. \\
\bottomrule
\end{tabular}
\end{table}

\paragraph{The supplied example can pass while another frozen check fails.}
In the original Angstrom conditions, 108 sessions received a regression example; 96 submitted patches passed it. Eighteen of those 96 failed another already-frozen security check. Eight were caught on another check of the same scenario; ten were caught only by variant checks. These 18 have not been shown to represent 18 independently qualified hidden behaviours. Dividing by all 108 would answer a different question: the probability of reaching that combination from the start of an assigned session. Conditioning on the 96 answers the maintainer who has already seen the test turn green. The result pertains to the other checks available for this one finding.

\paragraph{Even a passing battery can miss the reported defect.}
In the Ekubo study, 59 of 144 GPT-5.4-mini submissions passed both the original security checks and the ordinary-use suite. A nonblinded, model-led review of the apparent successes informed new checks, which exposed the audited defect in 17 of those 59. There was no independent human adjudication of these patches. The executable outcomes document missed cases in that collection. The later checks were informed by the patches they assessed. Applying the same proportion to future findings would overlook both that dependence and how the 59 were selected.

\paragraph{A stronger grader can also reach a ceiling.}
A later Ekubo study counted 48 Claude Code sessions with Opus; five interrupted starts were set aside and repeated. Half the counted sessions had no supplied regression test. All 48 passed the existing and follow-up checks mechanically. The configurations, exposure conditions and review procedures differ across the two Ekubo studies. We use the ceiling to motivate reporting which checks discriminate on which cases. The observations cannot rank the two models or establish an effect of supplying a particular example.

\section{Building the benchmark}
\label{sec:construction}
\paragraph{Finding population.}
The recruitment frame combines 140 historical task folders with a separate screen of public Code4rena and Sherlock final reports from 2022, 2023 and 2025. The archive is development material: some previously assessed families or findings could be familiar to the systems. The external screen used fixed publisher-by-year lists and a hash order within each list. Within a drawn report, the screen retained the publisher's final high- and medium-severity findings and recorded duplicates, exclusions and unresolved provenance. Cases reserved for a separate study were held out. Across 54 drawn reports, the publishers counted 646 high- or medium-severity findings. The screen preserved those entries and one additional finding explicitly marked out of scope, giving 647 records. Of these, 333 were selected for external follow-up in 28 provisional upstream families. A follow-up has enough documentary evidence to pursue. It has no automatic place in the benchmark.

To avoid collecting many revisions of the same component, the plan caps admissions at four per upstream family across both sources. Cases rotate across families before a later rank is considered. Within a reserve wave, an accessible landed revision tied to the finding and a contract test filename determine reading priority. All candidates in the frozen wave keep their place. Readers check shared code, aliases and overlapping fixes before treating two findings as independent. This samples a documented population of audited projects. It gives no estimate for Solidity projects in general.

\paragraph{Admitting one finding.}
The case record connects a public finding to an affected ordinary use, a compatible starting revision, and an attributable reference correction. The ordinary-use assessment must exercise the affected feature on both historical trees. A successful suite total is insufficient when the relevant tests were skipped or depend on an unavailable external service. Conversely, a logged failure elsewhere in a larger suite need not invalidate a separately identified passing use. Both facts stay in the record.

Example A must reveal the reported violation on the starting tree and support a permitted correction on the fixed tree. Readers inspect why each verdict occurred. A build failure, unrelated error, or expectation tied to an incidental implementation detail can produce a misleading change of verdict. The reference shows one acceptable correction. The grader must accept alternatives allowed by the requirement. Where a developer revision bundles unrelated work, the record documents attribution and compatibility or leaves the case pending.

The human admission reader checks eight claims covering identity and provenance, compatible revisions, the reported obligation, builds, ordinary use, A on each revision, and fair grading. A ninth decision admits the case against the full digest of a versioned evidence manifest. For one preselected case in every block of four, the human receives only the neutral source packet until their first judgment is locked or prior exposure is recorded. Assisted final decisions and independent first readings are recorded separately.

\paragraph{Qualifying another behaviour.}
The primary endpoint originally required a hidden observation on every finding. Before any expanded-benchmark repair session, a dated internal amendment removed that requirement. Readers had already seen historical repair results and evidence on seven new qualification cases. Of those seven only one supported a separately satisfiable hidden obligation. Four repeated the same observation, including variants with other inputs; one was dependent on the visible observation, and one measured ordinary use. Separately authored tests had often measured the same thing. Removing the requirement allows a single-obligation finding to contribute to RQ1 without silently broadening what a passing patch establishes.

Under the amended rule a curator reads the factual finding as well as the requirement. A second reader checks what the historical defect actually violated before H is assigned. Two observations qualify as distinct if there is a credible distinction in the caller-visible obligations: one can give a plausible correction of A that leaves H wrong and a plausible correction of H that leaves A wrong. A one-way diagnostic is recorded separately. A different input alone proves nothing about distinctness. The human signs H eligibility separately before repairs; the analysis also shows how conclusions change if evidence across two instances of the same mistake is omitted.

\section{Case-construction results}
Table~\ref{tab:intake} gives the external qualification counts as of 23 November 2026. The distinction between a scope recommendation and admission matters at this stage. The documentary reader cleared 38 of 160 reviewed findings, deferred 79 for a decision and 25 for additional evidence, and excluded 18. All 38 cleared cases have since reached a build attempt on both historical trees. An earlier snapshot had only 23 build attempts and 11 positive ordinary-use records.

\begin{table}[t]
\caption{External case qualification at the reporting cut. A case is counted once at each stage it reaches. Counts refer to findings, except where explicitly labelled. No row counts agent repairs.}
\label{tab:intake}
\small
\begin{tabular}{@{}p{.57\linewidth}rp{.27\linewidth}@{}}
\toprule
Stage reached & Count & Reference population \\
\midrule
Documentary follow-ups & 333 & 647 finding records \\
Completed benchmark scope reviews & 160 & 333 follow-ups \\
Cleared for technical qualification & 38 & 160 reviews \\
Both historical revisions built & 37 & 38 attempted \\
Specified ordinary use works offline on both & 18 & 38 attempted \\
Findings with a reviewed assessment for A & 9 & 18 offline cases \\
Assessment for A interpretable on both revisions & 8 & 9 reviewed \\
Versioned pages ready for human admission & 5 & 8 with validated A \\
Signed primary admissions & 0 & 5 prepared pages \\
\bottomrule
\end{tabular}
\end{table}

The main loss after compilation occurs at ordinary use. Seventy-four of 76 historical trees built. Only 18 of the 38 findings have positive evidence for affected ordinary use on both revisions offline. Some otherwise runnable tests need historical chain state or another external service. Access to that state would permit further assessment; it would not by itself establish ordinary-use success. Of the nine A assessments reviewed, eight have interpretable results on both revisions and one remains unresolved on the fixed revision. The eight span six upstream families. Five pages across four families are ready for the human reader. A technical pass has not been counted as a human decision.

Evidence for alternative examples is scarcer. The initial scope review cleared no paired cases among the 160. A separate intensive reading revisited twelve candidates across twelve families and judged that two had documentary support for the required distinctions. These are overlapping investigations with different depths of review. Neither of those two cases has signed paired eligibility. In the nine completed A assessments, the existing hidden-evidence review supports one distinct H. A more recent documentary comparison nominated one further candidate H across two instances. Nomination precedes assessment on both references and human sign-off.

\section{Prospective repair experiment}
\label{sec:experiment}
\paragraph{Units and allocation.}
The funded benchmark targets 80 admitted findings in at least 20 upstream families. The allocation includes 40 paired findings in at least 12 families. The configurations are Claude Code with Opus 5.5 and Codex CLI with GPT-5.4-mini, each at medium effort. They are complete repair systems with different tools as well as different models. The exact provider versions, repository images, resource limits and assignment seed are to be frozen before the first session.

Each configuration receives two independent sessions showing the standard example A for each admitted finding. On a paired finding the same configuration receives two additional sessions showing B. Thus the full design has \(80\times2\times2=320\) standard sessions and \(40\times2\times2=160\) additional sessions. The A sessions on a paired finding also contribute to the standard benchmark and are run only once. With twelve paired findings the corresponding count would be 368 and the randomized comparison would be exploratory. The smaller alternative requires a separate decision before seeing repair outcomes.

\paragraph{What changes between assignments.}
A paired case needs two independently justified behaviours from the same defect and a distinct common observation H that both groups will be graded on. The approved examples are designated A and B without consulting repair outcomes. Within each finding and configuration, two of four slots are randomly assigned A and two B. The remainder of the report, repository, ordinary-use suite, allowed tools, effort and time budget are identical within that block. A visible example cannot also occur in the common environment under a second name. Each session starts from the same pinned tree in a fresh workspace. The custodian retains the complete assignment and interleaves run order.

A scorer who cannot see the assignment grades both behaviours on both groups and also records H and ordinary use. A behavioural success requires a build, preserved ordinary use, and a reason-supported correction on that behaviour. The score permits a legitimate implementation different from the reference patch. A session ending in refusal, timeout or an empty patch remains in its assigned group. Infrastructure failures, unresolved assessments and retries retain their own records. First attempts stay visible throughout the analysis.

\paragraph{RQ1 and RQ2.}
For each configuration, let \(W\) indicate a working repair of the reported behaviour A. Average \(W\) over the two assigned A sessions for each finding, then give findings equal weight. Report the raw numerator and denominator and the number lost at patch production, build, ordinary use and A. Also weight families equally and show paired differences between configurations on the same findings. An unresolved assessment earns no demonstrated success in the primary rate; bounds also count unresolved cases at each extreme.

Let \(\mathcal{H}\) contain the primary cases with human-signed, distinct hidden evidence at the roster freeze. Among the same standard A attempts on \(\mathcal{H}\), report both \(W\) and \(W\) together with correction of H. Their difference measures the successes lost when that second obligation is added. Also show how many primary successes remain unevaluable on H, and how the cases in \(\mathcal{H}\) differ in family and reference route. Where H is missing inside the subset, its grade remains unresolved. Outside the subset no H grade is defined. A small difference on a narrow subset says little about other uses of the remaining contracts.

\paragraph{RQ3.}
For a paired finding and configuration, let \(Y_A\) and \(Y_B\) indicate working repairs of the respective behaviours on the same patch. For \(x,z\in\{A,B\}\), write \(\bar Y_{x\mid z}\) for the mean on behaviour \(x\) when example \(z\) was shown. The block contrast is
\begin{equation}
D = \tfrac{1}{2}\left[
  \big(\bar Y_{A\mid A}-\bar Y_{A\mid B}\big)
  +\big(\bar Y_{B\mid B}-\bar Y_{B\mid A}\big)\right].
\label{eq:contrast}
\end{equation}
For example, if every patch corrects A and none corrects B under either report, A is easier on these attempts but both differences are zero. If showing A helps A, while showing B makes little difference to B, the two terms reveal the asymmetry even when their mean is positive. Report both terms as well as their mean, averaged equally across finding-by-configuration blocks.

For the 40-pair design, the primary inference uses the random assignment of example labels. Permute the balanced labels within each block under the sharp no-effect null. Report an assignment-based 95\% interval for the average contrast and identify the large-sample approximation. Also report a bounded-score interval under independent block assignment; it governs the decision if the estimated standard error vanishes. With 80 independent finding-by-configuration blocks its conservative 95\% half-width is about 30 percentage points. With 24 blocks in the exploratory alternative it is about 55 points. These bounds describe the design and are not observed precision. Claim directional support only with a positive estimate, an interval excluding zero, a two-sided randomization \(p<.05\), and a positive lower bound after accounting for unresolved grades. Report the configurations separately, leave one family out at a time, and distinguish external findings from archive findings used in development. A wide interval leaves the question open. If nearly all patches satisfy both behaviours, report the remaining room for improvement under each assignment.

\paragraph{Results outstanding.}
There are no released sessions in the expanded benchmark, so RQ1 has no repair-rate denominator yet. RQ2 has no signed hidden subset, and RQ3 has no signed paired cases or observed assignment contrast. A future results section must show both the actual case flow and all assigned attempts, including unsuccessful and unresolved attempts. No performance estimate or significance test can be filled in from the qualification counts in Table~\ref{tab:intake}.

\section{What a result would let a maintainer conclude}

\label{sec:interpretation}

\subsection{Directional effect, inconclusive result or ceiling}

Suppose the funded comparison meets its directional decision rule. A maintainer could infer that, across the eligible cases and configurations under study, changing which example the agent received changed which specified behaviour its patch restored. The two component differences would show how that happened. One might account for most of the mean. A component could even be negative if that behaviour were repaired more often when the other example was shown. Reporting both directions and the results for each configuration would let the reader see whether the average describes either system well. A contrast below zero would call for reporting the reversal and its uncertainty.

A wide interval spanning zero would leave open effects that could matter to a maintainer. A narrow interval around zero would limit their plausible size on these cases. We have specified no equivalence margin. Where almost every patch repairs both behaviours under either assignment, we will report the successes, the shared failures, and how much room remained for improvement. The 48 counted Ekubo sessions in Table~\ref{tab:development} illustrate why that distinction matters. All reached the ceiling of the available checks. Deciding whether a caller can still obtain a wrong result would require evidence that can distinguish the alternatives.

\subsection{Limits, human judgment and conclusion}

The selection in Table~\ref{tab:qualification} favours findings with usable historical revisions and ordinary operations that can run offline. The requirements for H and paired admission narrow the eligible population further. Public findings or developer changes may also have been present in a model's training material. Withholding an assessment from the assigned session controls access during the study. It cannot establish what an agent encountered earlier. We will report results by family, weight families equally in a sensitivity analysis, omit one family at a time, and examine newly recruited families separately. Run order, provider versions and outages will help readers assess the assumption of independent session slots.

The first page ready for human judgment illustrates three limits on what we can call a working repair. The link between the finding and its reference fix was inferred during scope review. The audit report was the sole source for the requirement. The three named ordinary-use tests passed on both revisions while the complete project suite failed on both. Two named tests also assumed a particular allocation layout, although the behavioural requirement permitted others. A patch could meet that requirement and fail ordinary use as selected. We will report the behavioural grade and the named-test result separately so readers can see which decided the outcome. Three other possible endings of the private check were assessed by reading its source. None occurred in the stored runs. One reviewable page gives no estimate of admission speed or grading error on new patches.

Codex assisted with study design, interpretation of earlier runs, screening records, and manuscript preparation. Claude Code assisted with engineering and with preparing and reading case material. Model review informed the later Ekubo checks. In the qualification work a reader also helped identify the conditions it subsequently interpreted (Table~\ref{tab:historical-controls}). On the first reviewable page one model reader held five review roles. Agreement across those roles cannot count as five independent judgments. A fresh model session can check a version without having written it, but agreement between sessions still gives no estimate of human agreement. An assisted human admission decision records responsibility for a particular evidence version. Only a judgment recorded before exposure can contribute to an estimate of independent agreement. The scope records document early or incidental exposure in the later reviews, so treating every first recorded judgment as independent would overstate the available evidence. We plan an independent audit of 32 submitted patches by two human readers blind to assignment, system identity and automatic grades. No independent agreement rate is available at this cut. Manual assessment of functional effects in earlier repair work illustrates the purpose of such review.~\cite{bui2024apr4vul,bobadilla2026fixes}

The two development findings give a maintainer specific reasons to look past a passing example. At Angstrom another frozen check detected a problem among patches that passed the supplied test. At Ekubo later, patch-informed checks detected the audited defect among patches that had also passed ordinary-use checks. Technical qualification shows where the historical record can support a more precise question about a permitted outcome and where the evidence runs out. There is evidence that one case can reach human review. The recruitment yield and the first page leave the feasibility of the full roster open. Repair rates require signed cases and new submissions. Whether choosing the example changes the behaviour restored requires the randomized assignments and blind grades described here.


\section*{Data Availability}
The counts in this version come from archived development runs, dated screening tables, and versioned case records. No anonymous replication package has been deposited as of this draft. Case admission and the design freeze are still open. Early disclosure of protected judgments or withheld behaviours to study participants would compromise the comparison. For review, we plan an anonymized package with the screening rules, aggregate and redacted case tables, the analysis, the dated outcome amendment, and the evidence supporting each reported comparison. The final package will distinguish shareable upstream material from items for which only an identifier or a digest can be redistributed. Any material still sealed at submission will be identified with the reason and its release conditions.

\bibliographystyle{ACM-Reference-Format}
\bibliography{references}
\end{document}
